% --- ОБЩИЕ НАСТРОЙКИ ТЕКСТА ---
\onehalfspacing
\setlength{\parindent}{1.25cm} % Абзацный отступ

% Улучшенное выравнивание по ширине (меньше дырок и адекватные переносы)
\sloppy 
\hyphenpenalty=1000 
\tolerance=2000

% --- НАСТРОЙКА ОГЛАВЛЕНИЯ (СОДЕРЖАНИЕ) ---
\renewcaptionname{russian}{\contentsname}{\bfseries СОДЕРЖАНИЕ} % Жирный заголовок
\BeforeTOCHead[toc]{\renewcommand*{\raggedsection}{\centering}} % Центрирование заголовка

% Делаем пункты оглавления обычным шрифтом (не жирным)
\setkomafont{chapterentry}{\normalfont}
\setkomafont{chapterentrypagenumber}{\normalfont}
\KOMAoptions{toc=chapterentrydotfill} % Точки для глав

% --- НАСТРОЙКА ЗАГОЛОВКОВ (ОТСТУП 1.25 СМ) ---
\KOMAoptions{numbers=noendperiod} % Убирает точку в конце номера (1.1 вместо 1.1.)

% Устанавливаем шрифт 14pt жирный для всех заголовков
\setkomafont{disposition}{\rmfamily\bfseries\normalsize}

% Настройка отступов и интервалов (размеры beforeskip можно уменьшить, если заголовки кажутся слишком далеко)
\RedeclareSectionCommand[indent=1.25cm, beforeskip=12pt, afterskip=6pt]{chapter}
\RedeclareSectionCommand[indent=1.25cm, beforeskip=12pt, afterskip=6pt]{section}
\RedeclareSectionCommand[indent=1.25cm, beforeskip=12pt, afterskip=6pt]{subsection}

% --- СПИСКИ (enumitem) ---
\setlist{
	leftmargin=0pt,
	nosep,
	wide=\parindent,
	labelsep=0.5em,
	listparindent=\parindent
}
\setlist[itemize,1]{label=--} 
\setlist[enumerate,1]{label=\arabic*.}

% --- ПОДПИСИ И ТАБЛИЦЫ ---
\captionsetup{justification=centering, singlelinecheck=false}
\DeclareCaptionLabelFormat{simple}{#1~#2}
\captionsetup{labelformat=simple}

\DeclareCaptionLabelSeparator{emdash}{ --- }
\captionsetup[table]{
	labelsep=emdash,
	justification=raggedright,
	singlelinecheck=false,
	position=top,
	skip=6pt
}

% --- ЛИСТИНГИ КОДА ---
\lstset{
    basicstyle=\ttfamily\footnotesize,
    breaklines=true,
    texcl=true,
    frame=single
}

\DefineBibliographyStrings{russian}{
	urlseen = {дата обращения:},
}